\section{Continuum Description of Blood Flow}
\label{sec:continuum-description}
\label{sec:continuum-description-blood-flow}

This chapter fixes the continuum vocabulary used by the later reduced stenosis
model. The aim is not to repeat a full mechanics derivation. It is to record the
progression from material motion to transport identities, balance laws, stress
closure, and the incompressible flow equations from which cross-sectional
models inherit their variables and closure obligations.

\subsection{Continuum Scale and Field Variables}
\label{subsec:continuum-scale-and-fields}

At the large-vessel scale, blood-flow models replace individual blood cells,
platelets, and plasma constituents by macroscopic fields on a lumen. This is the
scale at which three-dimensional CFD, FSI, and one-dimensional area--flow models
can be placed in a common hierarchy. It is also a modeling assumption:
small-vessel, hematocrit-dependent, or low-shear regimes can require additional
rheology choices rather than the same constant-viscosity closure
\cite{KopplHelmig2023DimensionReducedModeling,GaldiEtAl2008HemodynamicalFlows,PriesEtAl1992BloodViscosityTubeFlow}.

\begin{assumption}[Large-vessel continuum setting]
    \label{ass:continuum-bridge-large-vessel-setting}
    Let $T_{\mathrm{final}}>0$,
    $I_T=(0,T_{\mathrm{final}})$, and
    $\ITBar=[0,T_{\mathrm{final}}]$. For each $t\in\ITBar$, the blood occupies
    a spatial lumen $\OmgBt\subset\R^3$, with space-time cylinder
    \begin{flalign*}
        \quad
        \mathcal Q_T=\{(t,\vx):t\in I_T,\ \vx\in\OmgBt\}. &&
    \end{flalign*}
    The continuum unknowns are the Eulerian velocity $\vu$, pressure $p$,
    density $\rho$, and Cauchy stress tensor $\vT$ on $\mathcal Q_T$.
\end{assumption}

These fields are the common language of the model hierarchy. A full-order
continuum tier evolves $\vu$ and $p$ on $\OmgBt$. A reduced 1D tier stores
section area, flow, and mean velocity obtained from these fields by additional
averaging and closure assumptions. The continuum assumption therefore fixes the
objects being reduced before any numerical discretization is introduced.

\begin{figure}[!htb]
    \centering
    \captionsetup{aboveskip=2pt,belowskip=2pt}
    \figtikz{compliant-vessel}
    \caption{Continuum-to-1D averaging notation for a compliant vessel. The
    physical lumen carries three-dimensional fields on $\OmgBt$; the reduced
    description records section averages such as area, flow, mean velocity, and
    gauge pressure. The figure is a bookkeeping map between model tiers, not a
    separate wall-traction or secondary-flow model.}
    \label{fig:compliant-vessel-reduction}
\end{figure}

\subsection{Flow Maps and Material Transport}
\label{subsec:continuum-flow-map-transport}

The two standard viewpoints are Eulerian and Lagrangian. The Eulerian viewpoint
observes fields at spatial locations $\vx$ in the current lumen. The Lagrangian
viewpoint follows a material label $\vxi$ from a reference configuration.

\begin{definition}[Flow-map vocabulary]
    \label{def:continuum-bridge-flow-map}
    A sufficiently regular flow map
    $\La_t:\OmgB(0)\to\OmgBt$ sends a material label $\vxi$ to its current
    position $\vx(t;\vxi)=\La_t(\vxi)$. The Eulerian velocity field generates
    the flow map through
    \begin{flalign*}
        \quad
        \pd_t\La_t(\vxi)=\vu(t,\La_t(\vxi)).
        &&
    \end{flalign*}
    For an Eulerian scalar field $\phi(t,\vx)$, the Lagrangian pullback is
    $\widehat{\phi}(t,\vxi)=\phi(t,\La_t(\vxi))$.
\end{definition}

The material derivative is the derivative of a field along this trajectory.
For a scalar field,
\begin{flalign*}
    \quad
    D_t\phi(t,\vx)
    =
    \dt\phi(t,\La_t(\vxi))
    =
    \pd_t\phi(t,\vx)+\vu(t,\vx)\cdot\nabla\phi(t,\vx),
    \qquad \vxi=\La_t^{-1}(\vx).
    &&
\end{flalign*}
For vector fields it is applied componentwise. This notation is the bridge
between material-volume statements and fixed-spatial differential equations.

\begin{theorem}[Reynolds transport bridge]
    \label{thm:continuum-bridge-reynolds-transport}
    Let $V(t)=\La_t(V(0))$ be a material volume. For a regular scalar field
    $\phi$,
    \begin{flalign*}
        \quad
        \dt\int_{V(t)}\phi\,\dx
        =
        \int_{V(t)}\left(D_t\phi+\phi\,\Div\vu\right)\dx
        =
        \int_{V(t)}\left(\pd_t\phi+\Div(\phi\vu)\right)\dx. &&
    \end{flalign*}
\end{theorem}

This identity supplies the modeling bridge used below: conservation statements
are first written on material volumes and then localized into Eulerian balance
laws.

\subsection{Mass, Momentum, and Stress Closure}
\label{subsec:continuum-balance-stress}

Mass conservation states that mass is neither created nor destroyed in material
volumes. In Eulerian form it is
\begin{flalign*}
    \quad
    \pd_t\rho+\Div(\rho\vu)=0
    \quad\Longleftrightarrow\quad
    D_t\rho+\rho\Div\vu=0. &&
\end{flalign*}
For constant density $\rho\equiv\rho_0>0$, this reduces to the incompressibility
constraint
\begin{flalign*}
    \quad
    \Div\vu=0. &&
\end{flalign*}

Linear momentum balance compares the rate of change of momentum in a material
volume with surface traction and body forces:
\begin{flalign*}
    \quad
    \dt\int_{V(t)}\rho\vu\,dV
    =
    \int_{\partial V(t)}\vT\nhat\,dS
    +
    \int_{V(t)}\rho\vfb\,dV. &&
\end{flalign*}
Using Reynolds transport, mass conservation, and localization gives
\begin{flalign*}
    \quad
    \rho\left(\pd_t\vu+(\vu\cdot\nabla)\vu\right)
    =
    \Divbf(\vT)+\rho\vfb. &&
\end{flalign*}
Here $\vfb$ is body force per unit mass.

The momentum equation becomes a closed flow model only after a constitutive
stress law is selected. With
\begin{flalign*}
    \quad
    \vct{D}(\vu)=\frac{1}{2}\left(\nabla\vu+\nabla\vu^\top\right),
    &&
\end{flalign*}
the Newtonian isotropic law is
\begin{flalign*}
    \quad
    \vT=-p\,\ident+2\eta\,\vct{D}(\vu)
    +\lambda\,\Div(\vu)\,\ident,
    &&
\end{flalign*}
where $\eta$ is dynamic viscosity and $\lambda$ is the volumetric-viscosity
coefficient. In the incompressible case this reduces to
\begin{flalign*}
    \quad
    \vT=-p\,\ident+2\eta\,\vct{D}(\vu). &&
\end{flalign*}
Generalized-Newtonian blood models keep the same stress structure but allow
viscosity to depend on a scalar shear-rate measure:
\begin{flalign*}
    \quad
    \dot\gamma=\sqrt{2\vct{D}(\vu):\vct{D}(\vu)},
    \qquad
    \vT=-p\,\ident
    +2\eta_{\mathrm{eff}}(\dot\gamma)\vct{D}(\vu).
    &&
\end{flalign*}
The Newtonian model is recovered when
$\eta_{\mathrm{eff}}$ is constant
\cite[Part~I, Ch.~1, Secs.~2.7--3.1]{GaldiEtAl2008HemodynamicalFlows}.
This large-vessel Newtonian baseline is a closure assumption for averaged
macroscopic behavior, not a microscopic statement about cells suspended in
plasma. The scale caveat matters: apparent viscosity can be nearly constant in
some large-vessel operating ranges, while diameter-, hematocrit-, and
shear-dependent behavior becomes more important in small-vessel or low-shear
settings \cite[Sec.~1.1.3, Fig.~1.14]{KopplHelmig2023DimensionReducedModeling}
\cite{PriesEtAl1992BloodViscosityTubeFlow}.

The displayed generalized-Newtonian form records the constitutive stress law
used in the review, not a theorem-level regularity statement for every scalar
viscosity curve. Positivity of
$\eta_{\mathrm{eff}}(\dot\gamma)$ alone is not used here as a standalone
parabolicity claim; such a claim would require stated lower-bound, growth, and
monotonicity hypotheses for
$2\eta_{\mathrm{eff}}(\sqrt{2D:D})D$.

\subsection{Incompressible Navier--Stokes Forms}
\label{subsec:continuum-navier-stokes}

Combining constant density, incompressibility, and Newtonian stress gives the
standard incompressible Navier--Stokes tier:
\begin{flalign*}
    \quad
    \begin{cases}
        \rho\left(\pd_t\vu+(\vu\cdot\nabla)\vu\right)
        =
        -\nabla p+\eta\Delta\vu+\rho\vfb, \\
        \Div\vu=0.
    \end{cases}
    &&
\end{flalign*}
Equivalently, with kinematic viscosity $\nu=\eta/\rho$,
\begin{flalign*}
    \quad
    \pd_t\vu+(\vu\cdot\nabla)\vu
    =
    -\rho^{-1}\nabla p+\nu\Delta\vu+\vfb,
    \qquad
    \Div\vu=0.
    &&
\end{flalign*}
The generalized-Newtonian incompressible form replaces the constant-viscosity
Laplacian by the divergence of the effective-viscosity stress:
\begin{flalign*}
    \quad
    \rho\left(\pd_t\vu+(\vu\cdot\nabla)\vu\right)
    =
    -\nabla p+
    \Divbf\!\left(2\eta_{\mathrm{eff}}(\dot\gamma)\vct{D}(\vu)\right)
    +\rho\vfb,
    \qquad
    \Div\vu=0.
    &&
\end{flalign*}
These equations are the full-order continuum reference for cardiovascular
models; reduced arterial models inherit their conservation structure after
geometric averaging and closure choices
\cite[Sec.~6]{QuarteroniFormaggia2004CardiovascularModeling}.

\subsection{Cross-Sectional Bridge to Area--Flow Models}
\label{subsec:continuum-area-flow-bridge}

The distributed 1D models reviewed later keep the axial conservation structure
of the continuum equations but replace the full three-dimensional field by
cross-sectional variables. Let $S(z,t)$ be the lumen section at axial coordinate
$z$, with outward-moving wall already included in the moving section. The
physical area and physical flow are
\begin{flalign*}
    \quad
    A(z,t)
    =
    \int_{S(z,t)}1\,dA,
    \qquad
    Q(z,t)
    =
    \int_{S(z,t)}u_z(z,r,\theta,t)\,dA,
    &&
\end{flalign*}
and the section mean velocity is $\bar u=Q/A$ when $A>0$. Applying mass
conservation to a short vessel slice and accounting for the wall motion gives
the reduced continuity law
\begin{flalign*}
    \quad
    \pd_t A+\pd_z Q=0. &&
\end{flalign*}

The axial momentum equation is not closed by $A$ and $Q$ alone. Its convective
term contains the second axial-velocity moment, commonly written through the
momentum-flux correction factor
\begin{flalign*}
    \quad
    \alpha(z,t)
    =
    \frac{A(z,t)}{Q(z,t)^2}
    \int_{S(z,t)}u_z(z,r,\theta,t)^2\,dA,
    &&
\end{flalign*}
with the usual limiting interpretation when $Q=0$. At review level, the
section-averaged axial balance has the structure
\begin{flalign*}
    \quad
    \pd_t Q
    +
    \pd_z\!\left(\alpha\frac{Q^2}{A}\right)
    +
    \frac{A}{\rho}\pd_z \bar p
    =
    \mathcal S_{\mathrm{wall}}
    +
    \mathcal S_{\mathrm{fric}}
    +
    \mathcal S_{\mathrm{geom}}
    +
    \mathcal S_{\mathrm{body}},
    &&
\end{flalign*}
where $\bar p$ is a section pressure representative. The reduced model is
therefore defined as much by its closures as by its conservation form:
$\alpha$ and the friction term encode the assumed velocity profile and rheology,
the pressure--area relation encodes wall mechanics, geometry terms record axial
variation of the reference radius or coordinate map, and inlet/outlet conditions
close the finite vessel problem
\cite{FormaggiaEtAl2003OneDimensionalBloodFlow,ShiEtAl2011BloodFlowModelReview,CanicEtAl2024Extended1DStenosis}.
Appendix~\ref{app:continuum-derivation-details} records the supporting
transport identities, while Section~\ref{sec:case-study} states the selected
solver-coordinate version used in the stenosis audit.

\subsection{Fixed and Moving Domains, Boundary Data, and Observables}
\label{subsec:continuum-domains-observables}

A fixed-domain blood-flow model treats the lumen geometry as prescribed. In the
simplest rigid-wall setting, $\OmgBt=\OmgB$ and the wall boundary carries a
no-slip or no-penetration condition while inlet and outlet portions carry
velocity, flow, pressure, traction, or reduced-network data. A moving-domain or
FSI setting instead evolves $\OmgBt$ through the wall motion and couples the
fluid velocity to the wall kinematics and stresses. Dimension-reduced models
replace parts of this boundary motion and stress information by section area,
wall-law, profile, friction, and rheology closures
\cite{KopplHelmig2023DimensionReducedModeling,QuarteroniFormaggia2004CardiovascularModeling}.

The principal observables depend on the selected tier and on the operator used
to read the field. Table~\ref{tab:continuum-principal-observables} records the
main categories used when a continuum field is compared with a reduced stenosis
model.

\begin{table}[!htb]
    \centering
    \small
    \caption{Principal continuum observables and the corresponding reduced
    quantities used in stenosis modeling.}
    \label{tab:continuum-principal-observables}
    \begin{tabularx}{\textwidth}{@{}
        >{\raggedright\arraybackslash}p{0.27\textwidth}
        >{\raggedright\arraybackslash}p{0.35\textwidth}
        >{\raggedright\arraybackslash}X
        @{}}
        \toprule
        \textbf{Continuum quantity}
            & \textbf{Typical observation}
            & \textbf{Reduced counterpart} \\
        \midrule
        Velocity $\vu$
            & Point, profile, section mean, or resolved field.
            & Mean axial velocity $\bar u=Q_{\mathrm{phys}}/A_{\mathrm{phys}}$
            and reconstructed profile descriptors. \\
        Pressure $p$
            & Local pressure, pressure drop, or pressure history.
            & Section-averaged pressure or wall-law gauge pressure. \\
        Flow through a section
            & $Q_{\mathrm{phys}}=\int_{S(z,t)} u_z\,dS$.
            & Stored flow state and physical-flow reconstruction. \\
        Wall traction and shear
            & Boundary stress or wall-shear measure.
            & Friction, profile, and wall closures rather than a resolved
            traction field. \\
        Current geometry
            & Lumen area, radius, and wall displacement.
            & Section area, reference geometry, and wall-law relation. \\
        \bottomrule
    \end{tabularx}
\end{table}

\begin{stenosisillustration}
The idealized stenosis case study inherits the continuum description in four
specific ways. First, the lumen is treated as a macroscopic blood domain
carrying velocity, pressure, density, and stress fields; blood cells are not
resolved individually. Second, the narrowing is represented through a declared
smooth vessel geometry and a selected wall treatment, so any wall motion not
solved at the continuum level must enter through the reduced wall law. Third,
mass and momentum balance are retained as the source of the section area and
flow equations, while discarded transverse velocity and stress information must
be supplied by profile, friction, rheology, and wall closures. Fourth, the
principal comparison quantities are continuum-derived observables: physical
flow, section-mean axial velocity, pressure when a pressure output is declared,
and velocity-profile descriptors.
\end{stenosisillustration}
